% Options for packages loaded elsewhere
\PassOptionsToPackage{unicode}{hyperref}
\PassOptionsToPackage{hyphens}{url}
\documentclass[
]{article}
\usepackage{xcolor}
\usepackage{amsmath,amssymb}
\setcounter{secnumdepth}{-\maxdimen} % remove section numbering
\usepackage{iftex}
\ifPDFTeX
  \usepackage[T1]{fontenc}
  \usepackage[utf8]{inputenc}
  \usepackage{textcomp} % provide euro and other symbols
\else % if luatex or xetex
  \usepackage{unicode-math} % this also loads fontspec
  \defaultfontfeatures{Scale=MatchLowercase}
  \defaultfontfeatures[\rmfamily]{Ligatures=TeX,Scale=1}
\fi
\usepackage{lmodern}
\ifPDFTeX\else
  % xetex/luatex font selection
\fi
% Use upquote if available, for straight quotes in verbatim environments
\IfFileExists{upquote.sty}{\usepackage{upquote}}{}
\IfFileExists{microtype.sty}{% use microtype if available
  \usepackage[]{microtype}
  \UseMicrotypeSet[protrusion]{basicmath} % disable protrusion for tt fonts
}{}
\makeatletter
\@ifundefined{KOMAClassName}{% if non-KOMA class
  \IfFileExists{parskip.sty}{%
    \usepackage{parskip}
  }{% else
    \setlength{\parindent}{0pt}
    \setlength{\parskip}{6pt plus 2pt minus 1pt}}
}{% if KOMA class
  \KOMAoptions{parskip=half}}
\makeatother
\usepackage{longtable,booktabs,array}
\usepackage{caption}
\captionsetup[table]{skip=6pt}
\newcounter{none} % for unnumbered tables
\usepackage{calc} % for calculating minipage widths
% Correct order of tables after \paragraph or \subparagraph
\usepackage{etoolbox}
\makeatletter
\patchcmd\longtable{\par}{\if@noskipsec\mbox{}\fi\par}{}{}
\makeatother
% Allow footnotes in longtable head/foot
\IfFileExists{footnotehyper.sty}{\usepackage{footnotehyper}}{\usepackage{footnote}}
\makesavenoteenv{longtable}
\setlength{\emergencystretch}{3em} % prevent overfull lines
\providecommand{\tightlist}{%
  \setlength{\itemsep}{0pt}\setlength{\parskip}{0pt}}
\usepackage{bookmark}
\IfFileExists{xurl.sty}{\usepackage{xurl}}{} % add URL line breaks if available
\urlstyle{same}
% fallback for those not using the hyperref driver hyperxmp:
\makeatletter
\@ifundefined{xmpquote}{\newcommand{\xmpquote}[1]{#1}}{}
\makeatother
\hypersetup{
  hidelinks,
  pdfcreator={LaTeX via pandoc}}

\author{}
\date{}

\begin{document}

\section{Paper IV v1.22-r3 --- Informe final consolidado de auditoría
externa adversarial
(E0--E8)}\label{paper-iv-v122-r3--informe-final-consolidado-de-auditoruxeda-externa-adversarial-e0e8}

\textbf{Ejecución:} \texttt{run\_v1.22\_r3}. Se completó el 2026-10-01,
con inicio a las 14:52:23 UTC y el E0 inicial a las 14:53:23 UTC; el E0
final está en la §11. La solicitó el propietario en el chat con el
mandato \texttt{EXTERNAL\_ADVERSARIAL\_REVALIDATION\_v1.22\_r3.md}.

\textbf{Veredicto global actual: PASS\_WITH\_FINDINGS.}

\begin{itemize}
\tightlist
\item
  Todas las comprobaciones del encargo se completaron.
\item
  NEW-04 (estado documental desactualizado) queda cerrado. Quedan
  cerradas también la acción de control de X-28 y el hallazgo E8-02.
\item
  Se añaden 12 filas semánticas nuevas, todas MATCH.
\item
  Queda abierto \textbf{un hallazgo MINOR nuevo, NEW-05}: dos términos
  españoles introducidos en los párrafos reescritos en r3. Exige una
  corrección de dos palabras antes de publicar.
\item
  No hay ningún defecto matemático, formal ni de identidad.
\item
  Siguen vigentes siete observaciones aceptadas, que son límites
  documentados y no exigen acción, más la mitigación X-28.
\item
  Esta revisión no sustituye una revisión por pares ni autoriza la
  publicación.
\end{itemize}

\begin{center}\rule{0.5\linewidth}{0.5pt}\end{center}

\subsection{1. Identidad actual}\label{1-identidad-actual}

{\def\LTcaptype{none} % do not increment counter
\begin{longtable}[]{@{}lll@{}}
\toprule\noalign{}
Objeto & Ruta & SHA-256 \\
\midrule\noalign{}
\endhead
\bottomrule\noalign{}
\endlastfoot
Objetivo r3 &
\texttt{02\_validation/02\_IA\_ADVERSARIAL\_AUDITS/AUDIT\_TARGET\_v1.22\_r3.json}
&
\texttt{df89c3dd33df19be5ddc4a42afaf5dd9053977623c7bc2a30fd12143b1044970}
(coincide con su sidecar) \\
Mandato r3 &
\texttt{…/EXTERNAL\_ADVERSARIAL\_REVALIDATION\_v1.22\_r3.md} &
\texttt{abfcaec9ae3ce5759b4177e0152bcd570e7ee31bb15e3f3df53553e6fb50d92f} \\
Manuscrito & \textbf{versión 1.22}, paquete \textbf{r3}, en
\texttt{01\_manuscript/v1.22\_editorial\_candidate\_r3/} & --- \\
Manifiesto r3 (226 archivos) &
\texttt{MANUSCRIPT\_REVIEW\_MANIFEST.json} &
\texttt{82217fe36719f4d2373bd63d30eb4efdb68e70b3ad95ef2f98217d3a6a5765bb} \\
ZIP r3 (228 miembros, CRC correcto) &
\texttt{PAPER\_IV\_v1.22\_REVIEW\_PACKAGE\_r3.zip} &
\texttt{c1acd8669d6cbcd89f7835933130416ae716c407250c5e690e5cf64fa837d2ec} \\
ES MD / TeX / PDF (73 p.) & \texttt{PAPER\_IV\_preprint\_v1.22\_es.*} &
\texttt{81db7bbb…} / \texttt{17c1306d…} / \texttt{d1048a62…} \\
EN MD / TeX / PDF (72 p.) & \texttt{PAPER\_IV\_preprint\_v1.22\_en.*} &
\texttt{d3415df3…} / \texttt{c6b14105…} / \texttt{4280b03e…} \\
Figuras (18 archivos) & \texttt{figures/}, \texttt{figures\_en/} &
idénticas a r2 \\
Freeze Lean (615 entradas) &
\texttt{05\_formalization/lean\_piv-v12-fb459343d234/} & manifiesto
\texttt{fb459343d234f968d7d32eff1491ea8a09aa2e135b313a80623012e7449042f5} \\
ZIP de fuentes Lean (615 miembros) &
\texttt{05\_formalization/LEAN\_SOURCE\_piv-v12-fb459343d234.zip} &
\texttt{cb2741454380736ffe01b59933057b5079fa9f865d5bbdd3f67534bf808a62e6} \\
Anexo (40 miembros) &
\texttt{05\_formalization/LEAN\_BOUNDED\_GAP\_ANNEX\_v1.0.zip} &
\texttt{2847a422865e06880d457ea806aae5b9f749d22359eff325349c09c01cb4837d} \\
Paquetes históricos &
\texttt{run\_v1.2\_r1/…/EXTERNAL\_AUDIT\_run\_v1.2\_r1.zip} &
\texttt{886ed7f0…f5e2b} \\
& \texttt{run\_v1.21\_r1/…/EXTERNAL\_REVALIDATION\_run\_v1.21\_r1.zip} &
\texttt{11659a97…7a638} \\
& \texttt{run\_v1.22\_r1/…/EXTERNAL\_REVALIDATION\_run\_v1.22\_r1.zip} &
\texttt{9c74d673…193a1fa} \\
& \texttt{run\_v1.22\_r2/…/EXTERNAL\_REVALIDATION\_run\_v1.22\_r2.zip} &
\texttt{626571be…1e480684} \\
\end{longtable}
}

Notas sobre la identidad:

\begin{itemize}
\tightlist
\item
  Los hashes completos están en
  \texttt{20\_EVIDENCE/E0/E0\_INITIAL.json}.
\item
  En r3 cambian el inglés y el español, como declara el editor.
  \textbf{No se exige ni se afirma identidad byte a byte del inglés con
  r2.} Lo que sí se exige y se verifica es la identidad del contenido
  protegido (E6), de Lean y de las figuras, y que los antecedentes no
  hayan cambiado.
\end{itemize}

\subsection{2. Veredictos actuales
E0--E8}\label{2-veredictos-actuales-e0e8}

{\def\LTcaptype{none} % do not increment counter
\begin{longtable}[]{@{}llllll@{}}
\toprule\noalign{}
Puerta & Veredicto actual & Alcance & Base & Procedencia & Evidencia \\
\midrule\noalign{}
\endhead
\bottomrule\noalign{}
\endlastfoot
E0 & \textbf{PASS} & identidad completa & \textbf{ejecutado en r3}
(inicial y final) & run\_v1.22\_r3 &
\texttt{20\_EVIDENCE/E0/E0\_INITIAL.json}, \texttt{E0\_FINAL.json} \\
E1 & \textbf{PASS} & 21 filas históricas + 12 nuevas & \textbf{heredado
y revalidado por identidad} (21) + \textbf{ejecutado en r3} (12) &
run\_v1.2\_r1 E1; run\_v1.22\_r3 &
\texttt{30\_REPORT/CLAIM\_MAP\_CONSOLIDATED.md};
\texttt{20\_EVIDENCE/E1/E1E3\_HEADERS\_CHECK.json} \\
E2 & \textbf{PASS} (observación R-01) & 56 componentes; las siete filas
de A.2 & \textbf{heredado y revalidado por identidad} & A2-1/2/7: v1.2;
A2-3/5/6: v1.21; A2-4: v1.22-r1 & \texttt{E2\_RECORD.md} de
run\_v1.2\_r1, run\_v1.21\_r1 y run\_v1.22\_r1 \\
E3 & \textbf{PASS} & cabeceras, identificadores, X-27 & \textbf{heredado
y revalidado por identidad} + cabeceras estáticas \textbf{ejecutadas en
r3} & run\_v1.2\_r1 → r3 & \texttt{20\_EVIDENCE/E3/E3\_RECORD.md} \\
E4 & \textbf{PASS --- reused\_verified\_external\_build} & 607 módulos,
19 targets, 224 export checks, 50 módulos del anexo & \textbf{heredado y
revalidado por identidad}; \textbf{sin build nuevo ni replay del kernel}
& build de run\_v1.2\_r1 (2026-09-30, 15:09--21:33 UTC) & consola
\texttt{ef5b01c7…}, records \texttt{0f2d8495…}, anexo
\texttt{57ce1895…}; \texttt{20\_EVIDENCE/E4/E4\_REUSE\_CHECK.json} \\
E5 & \textbf{PASS} & 59 comprobaciones + 7 controles & \textbf{heredado
y revalidado por identidad} & v1.2, v1.21, v1.22-r1 &
\texttt{20\_EVIDENCE/E5/E5\_RECORD.md} \\
E6 & \textbf{PASS\_WITH\_FINDINGS} (NEW-05) & delta EN+ES, 145 páginas,
PDF frente a TeX & \textbf{ejecutado en r3} & run\_v1.22\_r3 &
\texttt{20\_EVIDENCE/E6/E6\_RECORD.md}, \texttt{E6\_R3\_DELTA.json},
\texttt{PAGE\_INSPECTION\_LOG.csv} \\
E7 & \textbf{PASS} & bibliografía y atribución & \textbf{heredado y
revalidado por identidad} (referencias idénticas) & v1.2 → r2 &
\texttt{20\_EVIDENCE/E7/E7\_RECORD.md} \\
E8 & \textbf{PASS\_WITH\_OBSERVATIONS} (observación histórica X-28) &
verificador suplementario, 1266 logs, alcance 1266/844 &
\textbf{ejecutado en r3} + atribución heredada & run\_v1.21\_r1
(barrido); run\_v1.22\_r3 & \texttt{20\_EVIDENCE/E8/E8\_RECORD.md},
\texttt{E8\_R3\_CHECK.json} \\
\end{longtable}
}

La matriz por puerta (veredicto en cada revisión, comprobación nueva,
evidencia heredada y enlace de identidad) está en
\texttt{30\_REPORT/TRACEABILITY\_MATRIX.csv}.

\subsection{3. Historial original de
veredictos}\label{3-historial-original-de-veredictos}

Los veredictos INCONCLUSIVE se conservan como veredictos de sus
revisiones; no se reescriben.

{\def\LTcaptype{none} % do not increment counter
\begin{longtable}[]{@{}llll@{}}
\toprule\noalign{}
Revisión & Ejecución & Veredicto global original & Puertas no PASS \\
\midrule\noalign{}
\endhead
\bottomrule\noalign{}
\endlastfoot
v1.2 & run\_v1.2\_r1 (2026-09-30) & \textbf{INCONCLUSIVE} & E2
INCONCLUSIVE (A2-3, A2-4, A2-5, A2-6); E3/E6/E8 PASS\_WITH\_FINDINGS; E7
PASS\_WITH\_MINOR\_FINDINGS \\
v1.21 & run\_v1.21\_r1 & \textbf{INCONCLUSIVE} & E2 INCONCLUSIVE (A2-4);
E1 PASS\_WITH\_MINOR; E6 PASS\_WITH\_FINDINGS; E8
PASS\_WITH\_OBSERVATIONS \\
v1.22-r1 & run\_v1.22\_r1 & \textbf{PASS} con hallazgos MINOR no
bloqueantes & E6 PASS\_WITH\_FINDINGS; E8 PASS\_WITH\_OBSERVATIONS \\
v1.22-r2 & run\_v1.22\_r2 & \textbf{PASS\_WITH\_OBSERVATIONS} & E6, E8
PASS\_WITH\_OBSERVATIONS \\
v1.22-r3 & \textbf{run\_v1.22\_r3} & \textbf{PASS\_WITH\_FINDINGS} & E6
PASS\_WITH\_FINDINGS (NEW-05); E8 PASS\_WITH\_OBSERVATIONS \\
\end{longtable}
}

Todos los PASS históricos siguen siendo aplicables a r3, porque la
identidad está verificada:

\begin{itemize}
\tightlist
\item
  E1, E3, E5 y E7 de run\_v1.2\_r1;
\item
  el build de E4;
\item
  las aceptaciones de A.2 de v1.21 y v1.22-r1;
\item
  el cierre de X-15, NEW-01 y NEW-03 en r2.
\end{itemize}

\subsection{4. Mapa de afirmaciones y cobertura
matemática}\label{4-mapa-de-afirmaciones-y-cobertura-matemuxe1tica}

El mapa completo está en
\texttt{30\_REPORT/CLAIM\_MAP\_CONSOLIDATED.md}. Tiene dos tipos de
filas:

\begin{itemize}
\tightlist
\item
  \textbf{H}: las 21 filas históricas de run\_v1.2\_r1. Todas son MATCH
  salvo G.2, que en v1.2 quedó «no re-extraída» y ahora cubren N11 y
  N12.
\item
  \textbf{N}: 12 filas \textbf{nuevas de r3, diferenciadas}. No estaban
  en el E1 original.
\end{itemize}

Filas nuevas:

{\def\LTcaptype{none} % do not increment counter
\begin{longtable}[]{@{}llll@{}}
\toprule\noalign{}
Fila & Afirmación & Declaración & Veredicto \\
\midrule\noalign{}
\endhead
\bottomrule\noalign{}
\endlastfoot
N1 & Teorema 5.0, ventana reforzada &
\texttt{NearH1LocalConstructor.exists\_near\_partition\_paid\_by\_root\_sharp}
& MATCH \\
N2 & Teorema 5.0, interfaz anterior &
\texttt{…exists\_near\_partition\_paid\_by\_root} & MATCH \\
N3 & Corolario 5.4, geometría (5.17) &
\texttt{NearCriticalDichotomy.chordal\_far\_or\_criticalRoot} & MATCH.
Lean es algo más fuerte: añade
\textbar\textbar R\textbar−\textbar C\textbar\textbar{} ≤ n/100, que es
diferencia de tamaños, no diferencia simétrica. No se afirma que C sea
clique de G. \\
N4 & Corolario 5.4, densidad (5.18) &
\texttt{…chordal\_far\_or\_edge\_density} & MATCH, tomando el máximo de
ambos umbrales \\
N5 & Dicotomía (1.5) &
\texttt{HybridDichotomy.chordal\_far\_or\_nearStructure} & MATCH. Los
campos de \texttt{NearStructureWitness} contienen el comparador, una
raíz que es clique (\texttt{RegularizedRoot.isClique}), el
empaquetamiento del grafo original y las cuentas. \\
N6, N7 & Optimalidad del coeficiente cuadrático &
\texttt{SharpConstantOptimality.erdos81\_quadratic\_constant\_optimal},
\texttt{…\_isLeast} & MATCH. El caso real se obtiene por ampliación
racional, argumento válido. \\
N8 & Optimalidad del coeficiente lineal &
\texttt{LinearCoefficient.linear\_coefficient\_optimal} & MATCH.
Observación de tipo X-17: la cabecera es existencial en G. \\
N9 & Identidad (6.7b) &
\texttt{RootPartitionStability.sum\_rootPieceDefect\_eq} & MATCH \\
N10 & Proposición D.3 & \texttt{SplitMixedGap.mixed\_gap\_zero}
\textbf{y}
\texttt{SplitCompleteSharpValue.exists\_sharp\_cliquePartition\_allParities}
& MATCH. Hacen falta ambas declaraciones más (1.3a). \\
N11, N12 & G.2 / Proposición G.1 &
\texttt{FarExploration.CleanupRigidVerdict.threshold\_gt\_exp},
\texttt{…\_seventy\_two} & MATCH. Cierra el «no re-extraída» histórico
de G.2. \\
\end{longtable}
}

Cómo se verificaron las 13 cabeceras:

\begin{itemize}
\tightlist
\item
  Coinciden con los fuentes congelados (hash de manifiesto, línea, texto
  reextraído, namespace y variables en alcance) y con los logs del build
  externo (hash, módulo PASS, salida 0).
\item
  Son extracciones estáticas, \textbf{no tipos elaborados nuevos}.
\end{itemize}

\textbf{Las siete derivaciones de A.2}:

\begin{itemize}
\tightlist
\item
  A2-1 «§5.1, (5.4) y (5.7)»: ACCEPTABLE\_SUMMARY desde v1.2.
\item
  A2-2 «C.2, (3.8) y (3.10)»: ACCEPTABLE\_SUMMARY desde v1.2.
\item
  A2-3 «Lema 3.2 y C.1--C.2»: REQUIRES\_EXPANSION en v1.2;
  ACCEPTABLE\_SUMMARY desde v1.21.
\item
  A2-4 «C.3, Tabla C.1»: REQUIRES\_EXPANSION en v1.2 y v1.21;
  ACCEPTABLE\_SUMMARY desde v1.22-r1.
\item
  A2-5 «E.1»: ACCEPTABLE\_SUMMARY desde v1.21.
\item
  A2-6 «E.1, (E.2b), (E.2d)»: ACCEPTABLE\_SUMMARY desde v1.21.
\item
  A2-7 «D.1»: ACCEPTABLE\_SUMMARY desde v1.2.
\end{itemize}

\textbf{R-01} (cotas de segundo momento citadas en C.2) sigue vigente
como límite expositivo declarado.

\textbf{Límite de la revisión escrita:} es una comprobación a nivel de
resumen frente a enunciados formales congelados, no una revisión humana
por pares.

\subsection{5. E4: evidencia formal}\label{5-e4-evidencia-formal}

\textbf{Lo que se hizo} (run\_v1.2\_r1, intento 2):

\begin{itemize}
\tightlist
\item
  Runner propio del auditor, sin Lake. Lean 4.28.0, fuentes extraídas
  del ZIP (615/615) y objetos nuevos. Mathlib compartido en solo
  lectura, con la caché de 113 821 archivos sin cambios.
\end{itemize}

\textbf{Registros originales:}

\begin{itemize}
\tightlist
\item
  Consola: \texttt{modules\_planned\ 607,\ pass\ 607}, sin fallos, 607
  líneas \texttt{PASS\ n\ 607}, \texttt{EXIT\ 0}, y los 19 targets en
  PASS.
\item
  \texttt{records.jsonl}: 608 filas, todas PASS con salida 0.
\item
  \texttt{ReleaseExportCheck}: 224 \texttt{\#check} y 0 errores.
\item
  Anexo: 50/50 PASS, con 633 declaraciones y axiomas estándar.
\item
  Los 607 oleans son idénticos byte a byte a los del autor.
\end{itemize}

\textbf{Axiomas} (no son teoremas distintos):

\begin{itemize}
\tightlist
\item
  461 \emph{registros} de axiomas, de los cuales 314 son salidas de
  \texttt{\#print\ axioms}; el resto son repeticiones.
\item
  \texttt{\#print\ axioms} sobre \textbf{18 declaraciones principales}:
  siempre \{propext, Classical.choice, Quot.sound\}.
\item
  \textbf{Conos de 50 declaraciones públicas}: 0 axiomas no estándar, 0
  \texttt{sorryAx} y 0 \texttt{Erdos81}.
\end{itemize}

El \texttt{SUMMARY.json} de E4\_main, sobrescrito por AuditorChecks,
\textbf{no se usa}.

\textbf{Validez en r3:} r3 solo cambia prosa documental. E0 confirma que
Lean, el ZIP de fuentes, el anexo y la evidencia de run\_v1.2\_r1 no han
cambiado, y el script documental (sha \texttt{db2eddb8…}) da un
resultado idéntico al de r2. \textbf{No se ejecutó un build nuevo ni un
replay del kernel.}

\subsection{6. E5 y E8}\label{6-e5-y-e8}

\textbf{E5:}

\begin{itemize}
\tightlist
\item
  59 comprobaciones distintas: T01--T24, 21 de v1.21 y 14 de C.3. A
  ellas se suman 7 controles negativos añadidos, todos rechazados.
\item
  Las repeticiones no se suman como pruebas nuevas.
\item
  Las pruebas finitas y el certificado Certo (partición exacta de K3∨I3
  en 6 piezas, solo factibilidad) no sustituyen las pruebas universales.
\item
  Los controles del lector de logs pertenecen a E8 y no prueban
  teoremas.
\end{itemize}

\textbf{E8, el verificador suplementario
\texttt{verify\_frozen\_logs.py}:}

\begin{itemize}
\tightlist
\item
  Funciona con cierre seguro: un log ausente, sobrante, con hash
  distinto o con alerta lo hace fallar, igual que un recuento distinto
  de 1266.
\item
  Lo repetí en su sitio con salida propia: \textbf{PASS, 1266 logs}.
\item
  Lo repetí en una copia de prueba, cuyo único cambio es la ruta raíz:
  \textbf{15/15 casos unitarios y 4/4 corrupciones de CLI} con salida 1.
  El inventario y los casos coinciden con lo sellado, y el paquete no se
  sobrescribió.
\item
  \textbf{E8-02:} la regla de errores es literalmente la alternativa
  histórica de \texttt{audit\_publication.py}, sensible a mayúsculas,
  sin \texttt{re.I}. El refuerzo de sorry va en una regex separada.
  \textbf{Cerrado.}
\item
  \textbf{Cuatro agrupaciones:} 608 + 50 (auditor) + 422 + 186 (autor) =
  \textbf{1266 archivos, no 1266 pruebas}. Se cuentan una sola vez
  aunque se barrieron en v1.21, en r2 y dos veces en r3. El barrido
  antiguo del autor cubría 844 = 1266 − 422 (E8-01).
\item
  \textbf{Matiz histórico:} la regla congelada detecta el log negativo
  real solo por \texttt{sorryAx} en la impresión de axiomas, no por el
  aviso de \texttt{sorry}, que lleva backticks.
  \texttt{paperiv\_build.py} no lo detecta.
\end{itemize}

\textbf{X-28:}

\begin{itemize}
\tightlist
\item
  La acción de control de la entrega queda \textbf{cerrada}: el
  verificador es un control obligatorio del perfil de reproducción r3.
\item
  Se mantiene la \textbf{observación histórica}: los runners congelados
  conservan su patrón insuficiente y no se afirma que se hayan
  parcheado.
\item
  La evidencia de axiomas y conos (18 declaraciones y 50 conos) sigue
  siendo el control sustantivo.
\end{itemize}

\subsection{7. NEW-04 y paridad (E6)}\label{7-new-04-y-paridad-e6}

Hay cinco párrafos documentales por lengua: portada, §7, la nota de la
Tabla 7, A.2 y F.4. Los comparé íntegramente con la evidencia de r2.
\textbf{C.3 y las siete filas de A.2 ya no figuran como pendientes}, y
el estado de r2 se describe como PASS\_WITH\_OBSERVATIONS, sin inventar
un PASS limpio. Se conservan los INCONCLUSIVE históricos y los límites.
\textbf{NEW-04 queda CERRADO.}

Esos párrafos son registros fechados de r2. El dictamen sobre r3 es este
informe; no hace falta editar el manuscrito retroactivamente.

Comprobaciones:

\begin{itemize}
\tightlist
\item
  \textbf{Contenido protegido idéntico en EN y ES:} fórmulas en display
  y en línea, etiquetas, bloques de código, filas de tablas,
  encabezados, imágenes y bibliografía. Los dos códigos en línea
  añadidos son nombres de ejecuciones.
\item
  \textbf{Diffs:} el diff del editor coincide con el del auditor.
\item
  \textbf{El PDF corresponde a su TeX:} todas las regiones de diferencia
  en el texto extraído se explican por los párrafos cambiados.
\item
  \textbf{Páginas:} cambian EN 1, 33, 39, 40 y 64, y ES 1--6, 34--36,
  39, 40 y 66. Revisé visualmente las 145 páginas.
\end{itemize}

\textbf{NEW-05 (MINOR, nuevo).} El español reescrito en r3 usa dos
términos distintos de los ya establecidos en el manuscrito:

\begin{itemize}
\tightlist
\item
  «programación numérica de C.3» en A.2 (l.1355), cuando el texto usa
  «calendario numérico» y el encabezado de C.3 es «Un calendario
  explícito»;
\item
  «cota adaptada con muestra fijada» en F.4 (l.2356), cuando F.5 y la
  Tabla 10 dicen «muestras con anclajes».
\end{itemize}

\textbf{Acción antes de publicar:} restituir esos dos términos y repetir
el control de identidad ES y E6.

\subsection{8. Hallazgos: estado
actual}\label{8-hallazgos-estado-actual}

\texttt{FINDINGS.csv} tiene 53 filas: todas las históricas, más NEW-05 y
las correcciones del auditor C4-01, C4-02 y C4-03.

{\def\LTcaptype{none} % do not increment counter
\begin{longtable}[]{@{}lll@{}}
\toprule\noalign{}
Clase & IDs & Acción antes de publicar \\
\midrule\noalign{}
\endhead
\bottomrule\noalign{}
\endlastfoot
\textbf{Abierto} & NEW-05 (MINOR) & \textbf{Sí:} dos términos en ES \\
\textbf{Cerrados (29)} & X-01\ldots X-16, X-20, X-21, X-23, X-25, X-27,
NEW-01\ldots NEW-04, R-02, E8-01, E8-02, E8-03 & ninguna \\
\textbf{Mitigado} & X-28 & ninguna, siempre que el perfil de
reproducción mantenga obligatorio el verificador \\
\textbf{Observaciones aceptadas (7)} & X-17 (ahora también N8), X-18,
X-19, X-22, X-24, X-29, R-01 & ninguna; son límites documentados o
hechos históricos \\
\textbf{Información} & X-26 (independencia) & ninguna; se recomienda
revisión humana por pares antes de decidir publicar \\
\textbf{Correcciones conservadas (14)} & AUD-P1, AUD-C1, AUD-C2,
C2-01\ldots C2-08, C3-01\ldots C3-03, C4-01\ldots C4-03 & ninguna \\
\end{longtable}
}

Responsabilidades: el editor o autor corrigió los textos; el auditor
verificó cada cierre en la revisión indicada en \texttt{FINDINGS.csv}
(columna \texttt{revision\_verified\_r3}).

\subsection{9. Independencia y
límites}\label{9-independencia-y-luxedmites}

\begin{itemize}
\tightlist
\item
  \textbf{Modelo:} Claude Opus 5.5 (\texttt{claude-opus-5-5}), proveedor
  Anthropic.
\item
  \textbf{Sesión:} esta revalidación es \textbf{continuación de la misma
  sesión} de Claude Code que produjo las cuatro ejecuciones anteriores,
  con compactación de contexto.
\item
  \textbf{Independencia:} no es una sesión nueva, no es ciega y no es
  independiente de familia. El manuscrito declara el uso de Claude, que
  es la misma familia. Que haya continuidad entre revisiones no la
  convierte en independiente.
\item
  \textbf{Lecturas:} leí la evidencia previa y los documentos del
  editor, y los contrasté con evidencia propia. Todo está registrado en
  \texttt{00\_CONTROL/INPUT\_ACCESS\_LOG.csv}.
\item
  \textbf{Entorno:} misma máquina y misma caché que el build histórico.
\item
  \textbf{Ejecución:} no hubo ejecución dinámica nueva de Lean, ni
  build, ni replay del kernel. Se ejecutaron solo scripts Python de solo
  lectura y el verificador de logs del autor, en su sitio y en una
  copia.
\item
  \textbf{Inspección visual:} hojas 4-up a 70 dpi y comparaciones a 105
  dpi. Las imágenes que no pudieron confirmarse se volvieron a registrar
  desde hojas que sí se vieron (C4-01).
\end{itemize}

\subsection{10. Veredicto razonado}\label{10-veredicto-razonado}

\textbf{PASS\_WITH\_FINDINGS.}

\textbf{Por qué no es PASS limpio:} queda una acción abierta antes de
publicar, NEW-05. Es menor, pero real, y es una incoherencia
terminológica nueva introducida en r3. Por criterio del mandato no se
emite PASS con acciones abiertas.

\textbf{Por qué no es INCONCLUSIVE ni FAIL:}

\begin{itemize}
\tightlist
\item
  todas las comprobaciones se completaron;
\item
  no hay ningún defecto matemático, de correspondencia formal ni de
  identidad;
\item
  NEW-05 no altera ningún enunciado ni ninguna prueba.
\end{itemize}

\textbf{Qué certifica esta revisión:}

\begin{itemize}
\tightlist
\item
  la identidad exacta del paquete r3;
\item
  que el estado documental (NEW-04) es correcto y coherente con r2;
\item
  que el contenido protegido es idéntico a r2 y que cada PDF corresponde
  a su TeX;
\item
  que las 12 filas semánticas nuevas corresponden a sus declaraciones;
\item
  que el verificador suplementario funciona y protege el perfil de
  reproducción;
\item
  que los PASS históricos siguen enlazados por identidad a su evidencia
  original.
\end{itemize}

\textbf{Qué no certifica:}

\begin{itemize}
\tightlist
\item
  que el manuscrito esté listo para publicar mientras NEW-05 siga
  abierto;
\item
  una revisión humana por pares;
\item
  un build nuevo o un replay del kernel;
\item
  independencia de familia o de sesión;
\item
  novedad universal;
\item
  autorización para publicar.
\end{itemize}

\textbf{Si solo se corrige NEW-05} y una revalidación documental
confirma que únicamente cambiaron esas dos palabras en ES, no quedarían
acciones abiertas. El veredicto podría ser entonces PASS, manteniendo
explícitos los límites de la §9 y las siete observaciones aceptadas, que
son límites de alcance y no defectos.

\subsection{11. E0 final y sellado}\label{11-e0-final-y-sellado}

\begin{itemize}
\tightlist
\item
  \textbf{E0 final:} \texttt{20\_EVIDENCE/E0/E0\_FINAL.json}. El objeto
  \texttt{checks} es idéntico al inicial
  (\texttt{20\_EVIDENCE/logs/E0\_final.log}).
\item
  \textbf{Paquete:}
  \texttt{40\_PACKAGE/EXTERNAL\_REVALIDATION\_run\_v1.22\_r3.zip}, con
  su \texttt{RUN\_MANIFEST.json} y su sidecar SHA-256. Los ZIP
  históricos no se duplican; se enlazan por ruta y hash (§1).
\item
  \textbf{PDF:} generado con pandoc (Markdown → HTML y TeX
  independiente) más PyMuPDF Story (HTML → PDF). \textbf{No se compiló
  desde TeX.} Se revisó visualmente
  (\texttt{30\_REPORT/PDF\_VISUAL\_CHECK.csv}).
\item
  \textbf{Sin efectos externos:} no se publicó nada, no se hizo push ni
  se enviaron datos a terceros. No se modificaron el manuscrito, los
  originales ni los informes previos.
\end{itemize}

\end{document}
